\documentclass[a4paper,12pt]{scrartcl}

% ── Encoding & Language ─────────────────────────────────────────────────────
\usepackage[utf8]{inputenc}
\usepackage[T1]{fontenc}
\usepackage[english]{babel}

% ── Fonts ───────────────────────────────────────────────────────────────────
\usepackage{lmodern}

% ── Graphics & Colors ───────────────────────────────────────────────────────
\usepackage{xcolor}
\usepackage{graphicx}

% ── Tables ──────────────────────────────────────────────────────────────────
\usepackage{booktabs}
\usepackage{array}
\usepackage{longtable}
\usepackage{multirow}

% ── Code Listings ───────────────────────────────────────────────────────────
\usepackage{listings}

\definecolor{codebg}{RGB}{248,248,248}
\definecolor{codeframe}{RGB}{200,200,200}
\definecolor{kwcolor}{RGB}{0,0,180}
\definecolor{strcolor}{RGB}{163,21,21}
\definecolor{cmtcolor}{RGB}{0,128,0}
\definecolor{idcolor}{RGB}{30,80,120}

\lstdefinestyle{pythonstyle}{
  language=Python,
  backgroundcolor=\color{codebg},
  frame=single,
  rulecolor=\color{codeframe},
  basicstyle=\ttfamily\small,
  keywordstyle=\color{kwcolor}\bfseries,
  stringstyle=\color{strcolor},
  commentstyle=\color{cmtcolor}\itshape,
  identifierstyle=\color{idcolor},
  showstringspaces=false,
  breaklines=true,
  tabsize=4,
  numbers=left,
  numberstyle=\tiny\color{gray},
  numbersep=6pt,
  xleftmargin=18pt,
}

\lstdefinestyle{bashstyle}{
  language=bash,
  backgroundcolor=\color{codebg},
  frame=single,
  rulecolor=\color{codeframe},
  basicstyle=\ttfamily\small,
  keywordstyle=\color{kwcolor}\bfseries,
  stringstyle=\color{strcolor},
  commentstyle=\color{cmtcolor}\itshape,
  showstringspaces=false,
  breaklines=true,
  xleftmargin=10pt,
}

\lstdefinestyle{yamlstyle}{
  backgroundcolor=\color{codebg},
  frame=single,
  rulecolor=\color{codeframe},
  basicstyle=\ttfamily\small,
  showstringspaces=false,
  breaklines=true,
  xleftmargin=10pt,
  morecomment=[l]{\#},
  commentstyle=\color{cmtcolor}\itshape,
}

\lstdefinestyle{treestyle}{
  backgroundcolor=\color{codebg},
  frame=single,
  rulecolor=\color{codeframe},
  basicstyle=\ttfamily\footnotesize,
  showstringspaces=false,
  breaklines=false,
  xleftmargin=10pt,
}

% ── Links ────────────────────────────────────────────────────────────────────
\usepackage{hyperref}
\hypersetup{
  colorlinks=true,
  linkcolor=blue!70!black,
  urlcolor=blue!70!black,
  citecolor=blue!70!black,
  pdftitle={Relatorio GLSim},
  pdfauthor={},
}

% ── Portuguese captions (manual) ─────────────────────────────────────────────
\renewcommand{\contentsname}{Sum\'ario}
\renewcommand{\tablename}{Tabela}
\renewcommand{\figurename}{Figura}
\renewcommand{\lstlistingname}{C\'odigo}

% ── Spacing ──────────────────────────────────────────────────────────────────
\setlength{\parskip}{5pt}
\setlength{\parindent}{0pt}

% ── Compact lists ────────────────────────────────────────────────────────────
\newenvironment{compactitem}%
  {\begin{itemize}\setlength{\itemsep}{1pt}\setlength{\parsep}{0pt}}%
  {\end{itemize}}

\newenvironment{compactenum}%
  {\begin{enumerate}\setlength{\itemsep}{2pt}\setlength{\parsep}{0pt}}%
  {\end{enumerate}}

% ─────────────────────────────────────────────────────────────────────────────
\begin{document}

% ─── Capa ─────────────────────────────────────────────────────────────────────
\begin{titlepage}
  \centering
  \vspace*{2cm}
  {\Huge\bfseries Relat\'orio de An\'alise\par}
  \vspace{0.5cm}
  {\rule{\linewidth}{1.5pt}\par}
  \vspace{0.8cm}
  {\LARGE\bfseries GLSim\,---\,Global-Local Similarity\\[0.4em]
   para Reconhecimento Fino de Imagens\\[0.4em]
   com Vision Transformers\par}
  \vspace{0.8cm}
  {\rule{\linewidth}{1.5pt}\par}
  \vspace{1.5cm}

  \begin{tabular}{@{}rl@{}}
    \textbf{Artigo:}     & \href{https://arxiv.org/abs/2407.12891}{Global-Local Similarity for Efficient} \\
                         & \href{https://arxiv.org/abs/2407.12891}{Fine-Grained Image Recognition with ViT} \\[6pt]
    \textbf{Publicado:}  & ISCAS 2025 \\[6pt]
    \textbf{Autores:}    & Edwin Arkel Rios, Min-Chun Hu, Bo-Cheng Lai \\
                         & (NYCU\,/\,NTHU, Taiwan) \\[6pt]
    \textbf{Data:}       & 06/05/2026 \\
  \end{tabular}

  \vfill
  {\small Relat\'orio gerado a partir de an\'alise est\'atica do c\'odigo-fonte.}
\end{titlepage}

\tableofcontents
\clearpage

% ─── Seção 1 ──────────────────────────────────────────────────────────────────
\section{Vis\~ao Geral}

O GLSim \'e um framework de reconhecimento fino de imagens
(\textit{Fine-Grained Visual Recognition}\,---\,FGVR) baseado em Vision
Transformers (ViT). O problema central do FGVR \'e distinguir subclasses
visualmente muito parecidas (ex.: esp\'ecies de p\'assaros, modelos de carros,
cultivares de plantas), o que exige que o modelo foque em regi\~oes locais
discriminativas da imagem.

A contribui\c{c}\~ao principal \'e uma m\'etrica chamada \textbf{GLS
(Global-Local Similarity)}, computacionalmente barata, que identifica regi\~oes
discriminativas comparando a similaridade entre a representa\c{c}\~ao global da
imagem (token CLS) e as representa\c{c}\~oes locais (patches). Com base nisso,
a imagem \'e recortada, re-encodada e as \textit{features} s\~ao fundidas por
um m\'odulo agregador.

% ─── Seção 2 ──────────────────────────────────────────────────────────────────
\section{Estrutura de Diret\'orios}

\begin{lstlisting}[style=treestyle,
  caption={Estrutura de diret\'orios do reposit\'orio GLSim}]
GLSim-main/
|-- glsim/                     # Pacote principal
|   |-- model_utils/           # Arquiteturas de modelos
|   |   |-- glsim.py           # Modelo principal ViTGLSim
|   |   |-- configs.py         # Configuracoes de todas as variantes de ViT
|   |   |-- build_model.py     # Factory de modelos
|   |   |-- transformer.py     # Blocos Transformer customizados
|   |   |-- timm_glsim.py      # Integracao com timm (DINOv2, etc.)
|   |   |-- transfg.py         # Baseline: TransFG
|   |   |-- ffvt.py            # Baseline: FFVT
|   |   |-- cal.py             # Baseline: CAL
|   |   `-- resnet.py          # Baseline: ResNet
|   |-- data_utils/            # Pipeline de dados
|   |   |-- build_dataloaders.py
|   |   |-- build_transform.py
|   |   |-- datasets.py
|   |   `-- augmentations.py
|   |-- train_utils/           # Loop de treinamento
|   |   |-- trainer.py         # Classe Trainer principal
|   |   |-- scheduler.py       # Schedulers de learning rate
|   |   |-- focal_loss.py
|   |   |-- contrastive_loss.py
|   |   |-- mix.py             # CutMix / MixUp
|   |   `-- scaler.py          # AMP (mixed precision)
|   `-- other_utils/
|       |-- build_args.py      # Parsing de argumentos CLI
|       `-- yaml_config_hook.py
|-- tools/                     # Scripts executaveis
|   |-- train.py               # Treinamento e avaliacao
|   |-- inference.py           # Inferencia em imagens individuais
|   |-- vis_dfsm.py            # Visualizacao do mecanismo discriminativo
|   |-- heatmap.py
|   |-- calc_flops.py
|   |-- compute_feature_metrics.py
|   |-- demo.py                # Demo Gradio
|   `-- preprocess/            # Scripts de preparacao de datasets
|-- configs/                   # Arquivos YAML de configuracao
|   |-- methods/glsim.yaml     # Config do metodo GLSim
|   |-- datasets/*.yaml        # Config de cada dataset
|   |-- augs/*.yaml            # Config de augmentacoes
|   `-- settings/ft_is224.yaml # Config base de fine-tuning
|-- data/                      # CSVs de splits (train/val/test)
|-- samples/                   # Imagens de exemplo
|-- assets/                    # Figuras do paper
|-- requirements.txt
`-- setup.py
\end{lstlisting}

% ─── Seção 3 ──────────────────────────────────────────────────────────────────
\section{Arquitetura do Modelo}

\subsection{Componentes Principais
  (\texttt{glsim/model\_utils/glsim.py})}

A classe central \'e \texttt{ViTGLSim(nn.Module)}, composta por quatro
m\'odulos em sequ\^encia:

\subsubsection{(a) Patch Embedding + Tokeniza\c{c}\~ao}

\begin{lstlisting}[style=pythonstyle,
  caption={Defini\c{c}\~ao do patch embedding}]
self.patch_embedding = nn.Conv2d(in_channels=3, out_channels=768,
                                  kernel_size=(16, 16), stride=(16, 16))
\end{lstlisting}

\begin{compactitem}
  \item Imagem \texttt{(B, 3, 224, 224)} $\rightarrow$ patches
        \texttt{(B, 196, 768)} (para ViT-B/16)
  \item Token CLS aprendiv\'el concatenado: sequ\^encia fica
        \texttt{(B, 197, 768)}
  \item \textit{Positional embedding} somado
\end{compactitem}

\subsubsection{(b) Encoder Transformer}

Implementado em \texttt{glsim/model\_utils/transformer.py}:

\begin{compactitem}
  \item \texttt{num\_layers=12}, \texttt{hidden\_size=768},
        \texttt{num\_heads=12}, \texttt{ff\_dim=3072} (ViT-B/16)
  \item Cada bloco: LayerNorm $\to$ Multi-head Self-Attention
        $\to$ Residual $\to$ LayerNorm $\to$ FFN $\to$ Residual
  \item Suporte a \textbf{Stochastic Depth (DropPath)}
        para regulariza\c{c}\~ao
  \item Retorna \textit{features} intermedi\'arias e mapas de aten\c{c}\~ao
        para visualiza\c{c}\~ao
\end{compactitem}

\subsubsection{(c) M\'odulo GLSCM --- Global-Local Similarity Crop Module}

\begin{lstlisting}[style=pythonstyle,
  caption={L\'ogica do GLSCM (pseudoc\'odigo)}]
get_crops(x, images):
    g = x[:, 0, :]      # token CLS (rep. global), shape: (B, 1, 768)
    l = x[:, 1:, :]     # patches restantes (rep. locais), shape: (B, 196, 768)
    sim = cosine_similarity(g, l)   # similaridade por patch: (B, 196)
    # top-k patches mais similares -> bounding box 2D -> recorte + upsample
\end{lstlisting}

\textbf{Intui\c{c}\~ao:} Os patches mais similares ao CLS s\~ao os mais
representativos do conte\'udo global\,---\,logo, as regi\~oes discriminativas
da subclasse. O recorte dessas regi\~oes \'e redimensionado para o tamanho
original e re-encodado.

\textbf{Alternativas de m\'etrica de similaridade} (configur\'avel via
\texttt{sim\_metric}):
\begin{compactitem}
  \item \texttt{cos} (padr\~ao): similaridade de cosseno, m\'aximos
        selecionados
  \item \texttt{l1} / \texttt{l2}: dist\^ancia, m\'inimos selecionados
\end{compactitem}

\subsubsection{(d) Aggregator}

\begin{compactitem}
  \item Transformer leve com
        \texttt{aggregator\_num\_hidden\_layers=1} (padr\~ao)
  \item Entrada: \texttt{[CLS\_original, CLS\_crop]}
        $\to$ shape \texttt{(B, 2, 768)}
  \item Combina as duas representa\c{c}\~oes para a predi\c{c}\~ao final
  \item Seguido de LayerNorm
\end{compactitem}

\subsubsection{(e) Cabe\c{c}a de Classifica\c{c}\~ao}

\texttt{nn.Linear(768, num\_classes)} aplicado no token CLS do aggregator.

% ── Fluxo de dados ───────────────────────────────────────────────────────────
\subsection{Fluxo de Dados Completo}

\begin{lstlisting}[style=treestyle,
  caption={Pipeline completo de dados do GLSim}]
Imagem (B, 3, 224, 224)
        |
        v  patchify_tokenize
  Tokens (B, 197, 768)
        |
        v  forward_encoder (ViT)
  Features (B, 197, 768)
        |
   +---------+
   |         |  get_crops -> recorte (B, 3, 224, 224)
   |         v  patchify_tokenize + forward_encoder
   |   Features Crop (B, 197, 768)
   |         |
   +---->cat (B, 394, 768)
        |
        v  forward_reducer (extrai CLS de cada metade)
  [CLS_orig, CLS_crop] (B, 2, 768)
        |
        v  aggregator (Transformer 1 camada)
  Features agregadas (B, 2, 768)
        |
        v  head (Linear)
  Logits (B, num_classes)
\end{lstlisting}

\subsection{Variantes de Backbone Suportadas}

\begin{table}[h!]
  \centering
  \caption{Variantes de backbone ViT dispon\'iveis no GLSim}
  \begin{tabular}{lcccc}
    \toprule
    \textbf{Variante} & \textbf{hidden\_size} & \textbf{heads} &
    \textbf{layers} & \textbf{patch} \\
    \midrule
    ViT-T/16             & 192  & 3  & 12 & $16{\times}16$ \\
    ViT-S/16             & 384  & 6  & 12 & $16{\times}16$ \\
    ViT-B/16             & 768  & 12 & 12 & $16{\times}16$ \\
    ViT-L/16             & 1024 & 16 & 24 & $16{\times}16$ \\
    ViT-H/14             & 1280 & 16 & 32 & $14{\times}14$ \\
    DINOv2 (via timm)    & 768  & 12 & 12 & $14{\times}14$ \\
    \bottomrule
  \end{tabular}
  \label{tab:backbones}
\end{table}

% ─── Seção 4 ──────────────────────────────────────────────────────────────────
\section{Configura\c{c}\~ao do M\'etodo GLSim}

Arquivo: \texttt{configs/methods/glsim.yaml}

\begin{lstlisting}[style=yamlstyle,
  caption={Configura\c{c}\~ao principal do m\'etodo GLSim}]
classifier: 'cls'         # usa token CLS como representacao
dynamic_anchor: True      # tamanho do recorte = tamanho de um patch
anchor_class_token: True  # token CLS separado para o recorte
reducer: 'cls'            # extrai apenas o token CLS de cada stream
aggregator: True          # habilita o modulo aggregator
aggregator_norm: True     # LayerNorm apos aggregator
\end{lstlisting}

O modo \texttt{dynamic\_anchor=True} \'e a configura\c{c}\~ao principal: o
\textit{bounding box} do recorte \'e determinado automaticamente pelo GLSCM a
cada \textit{forward pass}, adaptando-se ao conte\'udo da imagem.

% ─── Seção 5 ──────────────────────────────────────────────────────────────────
\section{Datasets Suportados}

\begin{table}[h!]
  \centering
  \caption{Datasets suportados pelo GLSim}
  \begin{tabular}{lll}
    \toprule
    \textbf{Dataset}   & \textbf{Dom\'inio}            & \textbf{Classes} \\
    \midrule
    CUB-200-2011       & Aves                          & 200        \\
    Stanford Cars      & Ve\'iculos                    & 196        \\
    FGVC Aircraft      & Aeronaves                     & 100        \\
    NABirds            & Aves (Am\'erica do Norte)     & 555        \\
    iNat17             & Natureza geral                & 5.089      \\
    DAFB               & Personagens anime             & $\sim$2.000 \\
    Oxford Pets        & Animais dom\'esticos          & 37         \\
    Oxford Flowers     & Flores                        & 102        \\
    Stanford Dogs      & C\~aes                        & 120        \\
    Food-101           & Alimentos                     & 101        \\
    Cotton / Soy       & Plantas agr\'icolas           & varia      \\
    VegFru             & Vegetais/frutas               & varia      \\
    ImageNet           & Geral                         & 1.000      \\
    \bottomrule
  \end{tabular}
  \label{tab:datasets}
\end{table}

% ─── Seção 6 ──────────────────────────────────────────────────────────────────
\section{Pipeline de Treinamento}

\subsection{Augmenta\c{c}\~oes (\texttt{configs/augs/})}

\begin{compactitem}
  \item \textbf{weakaugs}: flip horizontal, normaliza\c{c}\~ao padr\~ao
  \item \textbf{medaugs}: flip + RandAugment leve
  \item \textbf{strongaugs}: TrivialAugmentWide, Random Erasing
\end{compactitem}

\subsection{Otimizador e Scheduler}

\begin{compactitem}
  \item Otimizador: SGD (via \texttt{timm.optim.create\_optimizer})
  \item Scheduler: cosine decay com warmup
        (\texttt{glsim/train\_utils/scheduler.py})
  \item Learning rate padr\~ao: \texttt{0.01}
  \item Suporte a AMP (\texttt{fp16}) via \texttt{torch.cuda.amp.autocast}
        + \texttt{NativeScaler}
\end{compactitem}

\subsection{Fun\c{c}\~oes de Perda Dispon\'iveis}

\begin{compactitem}
  \item \texttt{CrossEntropyLoss} (padr\~ao)
  \item \texttt{LabelSmoothingCrossEntropy} (\texttt{-{}-ls})
  \item \texttt{FocalLoss} (\texttt{-{}-focal\_gamma})
  \item Contrastive Loss (\texttt{contrastive\_loss.py}) para TransFG
\end{compactitem}

\subsection{T\'ecnicas de Regulariza\c{c}\~ao}

\begin{compactitem}
  \item Label Smoothing
  \item CutMix\,/\,MixUp (\texttt{-{}-cm}, \texttt{-{}-mu})
  \item Stochastic Depth (\texttt{-{}-sd})
  \item Gradient Accumulation
        (\texttt{-{}-gradient\_accumulation\_steps})
  \item Gradient Clipping (\texttt{-{}-clip\_grad})
\end{compactitem}

\subsection{Treinamento Distribu\'ido}

\begin{compactitem}
  \item Suporte a \texttt{DistributedDataParallel} (DDP)
  \item \texttt{RASampler} para Repeated Augmentation em treino
        distribu\'ido
\end{compactitem}

\subsection{Experimentos e Monitoramento}

\begin{compactitem}
  \item Integra\c{c}\~ao nativa com \textbf{Weights \& Biases (WandB)}
  \item Salva modelo \textit{best}, \textit{last} e por \'epoca
        configur\'avel
  \item Log de acur\'acia top-1 e top-5, loss, LR a cada
        \texttt{log\_freq} itera\c{c}\~oes
\end{compactitem}

% ─── Seção 7 ──────────────────────────────────────────────────────────────────
\section{Baselines Inclu\'idos para Compara\c{c}\~ao}

\begin{table}[h!]
  \centering
  \caption{Modelos baseline dispon\'iveis no reposit\'orio}
  \begin{tabular}{lll}
    \toprule
    \textbf{Modelo}         & \textbf{Arquivo}            & \textbf{Refer\^encia} \\
    \midrule
    TransFG                 & \texttt{transfg.py}         & Trans Fine-Grained \\
    FFVT                    & \texttt{ffvt.py}            & Feature Fusion ViT \\
    CAL                     & \texttt{cal.py}             & Counterfactual Attention Learning \\
    ResNet                  & \texttt{resnet.py}          & ResNet-50/101 \\
    ViT puro (sem GLSim)    & \texttt{build\_model.py}    & --- \\
    \bottomrule
  \end{tabular}
  \label{tab:baselines}
\end{table}

% ─── Seção 8 ──────────────────────────────────────────────────────────────────
\section{Ferramentas Utilit\'arias (\texttt{tools/})}

\begin{table}[h!]
  \centering
  \caption{Scripts utilit\'arios em \texttt{tools/}}
  \begin{tabular}{>{\ttfamily}lp{9cm}}
    \toprule
    \normalfont\textbf{Script} & \textbf{Fun\c{c}\~ao} \\
    \midrule
    train.py                      & Treinamento completo e avalia\c{c}\~ao no test set \\[2pt]
    inference.py                  & Infer\^encia em imagem(ns) \'unica(s); salva original + recorte lado a lado \\[2pt]
    vis\_dfsm.py                  & Visualiza o mecanismo discriminativo (GLS vs.\ attention rollout) em lote \\[2pt]
    heatmap.py                    & Gera heatmaps de aten\c{c}\~ao \\[2pt]
    calc\_flops.py                & Calcula FLOPs e par\^ametros do modelo \\[2pt]
    compute\_feature\_metrics.py  & M\'etricas de features (separabilidade, etc.) \\[2pt]
    demo.py                       & Interface Gradio interativa \\[2pt]
    preprocess/                   & Scripts para preparar splits CSV de cada dataset \\
    \bottomrule
  \end{tabular}
  \label{tab:tools}
\end{table}

% ─── Seção 9 ──────────────────────────────────────────────────────────────────
\section{Depend\^encias Principais}

\begin{table}[h!]
  \centering
  \caption{Principais bibliotecas utilizadas pelo GLSim}
  \begin{tabular}{lll}
    \toprule
    \textbf{Biblioteca} & \textbf{Vers\~ao} & \textbf{Uso} \\
    \midrule
    PyTorch  & ---    & framework base \\
    timm     & 0.9.12 & modelos pr\'e-treinados, otimizadores, losses \\
    einops   & ---    & opera\c{c}\~oes tensoriais leg\'iveis \\
    wandb    & ---    & rastreamento de experimentos \\
    gradio   & ---    & demo interativa \\
    scipy    & 1.8.0  & utilit\'arios cient\'ificos \\
    \bottomrule
  \end{tabular}
  \label{tab:deps}
\end{table}

% ─── Seção 10 ─────────────────────────────────────────────────────────────────
\section{Uso como M\'odulo Python}

\begin{lstlisting}[style=pythonstyle,
  caption={Exemplo de uso do \texttt{ViTGLSim} como m\'odulo Python}]
import torch
from glsim.model_utils import ViTGLSim, ViTConfig

cfg = ViTConfig(
    model_name='vit_b16',
    debugging=True,
    classifier='cls',
    dynamic_anchor=True,
    reducer='cls',
    aggregator=True,
    aggregator_norm=True
)
model = ViTGLSim(cfg)

x = torch.rand(2, 3, 224, 224)   # batch de 2 imagens RGB 224x224
out = model(x)  # retorna (logits, crops) quando anchor_size ativo
\end{lstlisting}

% ─── Seção 11 ─────────────────────────────────────────────────────────────────
\section{Comandos Principais}

\subsection{Instala\c{c}\~ao}

\begin{lstlisting}[style=bashstyle, caption={Instala\c{c}\~ao do pacote}]
pip install -e .
\end{lstlisting}

\subsection{Treinamento (CUB, ViT-B/16, 224px)}

\begin{lstlisting}[style=bashstyle,
  caption={Treinamento no dataset CUB-200-2011 com ViT-B/16}]
python tools/train.py \
  --cfg configs/cub_ft_is224_medaugs.yaml \
  --lr 0.01 \
  --model_name vit_b16 \
  --cfg_method configs/methods/glsim.yaml
\end{lstlisting}

\subsection{Avalia\c{c}\~ao de Checkpoint}

\begin{lstlisting}[style=bashstyle,
  caption={Avalia\c{c}\~ao de um checkpoint salvo}]
python tools/train.py --ckpt_path ckpts/cub_glsim_224.pth --test_only
\end{lstlisting}

\subsection{Infer\^encia em Imagem \'Unica}

\begin{lstlisting}[style=bashstyle,
  caption={Infer\^encia com visualiza\c{c}\~ao do recorte discriminativo}]
python tools/inference.py \
  --ckpt_path ckpts/dafb_glsim.pth \
  --images_path samples/others/dafb_rena_170785.jpg \
  --vis_mask glsim_norm
\end{lstlisting}

\subsection{Visualiza\c{c}\~ao do Mecanismo Discriminativo}

\begin{lstlisting}[style=bashstyle,
  caption={Visualiza\c{c}\~ao do GLSCM em lote (NABirds + DINOv2)}]
python tools/vis_dfsm.py \
  --serial 52 --batch_size 4 --vis_cols 4 \
  --cfg configs/nabirds_ft_is224_weakaugs.yaml \
  --model_name glsvit_base_patch14_dinov2.lvd142m \
  --vis_mask gls --vis_mask_pow
\end{lstlisting}

% ─── Seção 12 ─────────────────────────────────────────────────────────────────
\section{Resultados Reportados no Paper}

O m\'etodo obt\'em resultados competitivos em m\'ultiplos benchmarks de FGVR.
Os destaques s\~ao:

\begin{compactitem}
  \item \textbf{CUB-200-2011:} Acur\'acia top-1 compar\'avel ao estado da
        arte em 224px e 448px.
  \item \textbf{NABirds\,/\,iNat17:} Vantagem em datasets de grande escala.
  \item \textbf{Efici\^encia:} O GLSCM tem custo computacional ordens de
        magnitude menor que \textit{attention rollout}, com FLOPs totais
        similares ao ViT base.
  \item Checkpoints pr\'e-treinados dispon\'iveis no HuggingFace:
        \href{https://huggingface.co/NYCU-PCSxNTHU-MIS/GLSim}%
             {\texttt{NYCU-PCSxNTHU-MIS/GLSim}}.
\end{compactitem}

% ─── Seção 13 ─────────────────────────────────────────────────────────────────
\section{Pontos T\'ecnicos de Destaque}

\begin{compactenum}
  \item \textbf{Sem necessidade de anota\c{c}\~oes de localiza\c{c}\~ao:}
        O GLSCM detecta regi\~oes discriminativas automaticamente, sem
        \textit{bounding boxes} supervisionados.

  \item \textbf{Custo baixo do GLSCM:} A similaridade cosseno entre o CLS
        e os patches \'e uma opera\c{c}\~ao $\mathcal{O}(N)$ sobre a
        sequ\^encia de patches\,---\,muito mais barata que
        \textit{attention rollout} ($\mathcal{O}(N^2)$ por camada).

  \item \textbf{Dynamic anchor:} O tamanho do recorte \'e determinado
        dinamicamente pelos \'indices dos top-$k$ patches, adaptando-se
        ao tamanho do objeto em cada imagem.

  \item \textbf{Modularidade:} O c\'odigo separa claramente backbone,
        mecanismo de crop, aggregator e head. \'E poss\'ivel usar o
        \texttt{ViTGLSim} diretamente como m\'odulo ou via
        \texttt{build\_model} com qualquer backbone do timm.

  \item \textbf{Test-time augmentation:} Suporte a flip horizontal no
        momento de avalia\c{c}\~ao (\texttt{test\_flip}), que melhora
        marginalmente a acur\'acia.
\end{compactenum}

% ─── Seção 14 ─────────────────────────────────────────────────────────────────
\section{Limita\c{c}\~oes Identificadas no C\'odigo}

\begin{compactitem}
  \item \texttt{F.upsample\_bilinear} est\'a deprecado no PyTorch moderno
        (deveria ser \texttt{F.interpolate(...,\,mode=`bilinear')}).

  \item O trainer usa \texttt{torch.cuda.synchronize()} no loop de treino
        mesmo sem necessidade em modo n\~ao-distribu\'ido, adicionando
        overhead desnecess\'ario.

  \item A fun\c{c}\~ao \texttt{imshow} para visualiza\c{c}\~ao de erros
        chama \texttt{input()} interativo, incompat\'ivel com ambientes
        headless/servidor.

  \item \texttt{setup.py} usa o padr\~ao legado em vez de
        \texttt{pyproject.toml}.
\end{compactitem}

\end{document}
